\documentclass[../main.tex]{subfiles}
\begin{document}
%==========================================TIÊU ĐỀ TẬP========================================
\begin{table}[h]
    \begin{adjustwidth}{-1cm}{}
        \begin{tabular}{>{\centering\arraybackslash}p{6cm}|>{\raggedright\arraybackslash}p{14cm}}
            \multirow{5}{6cm}
            % Cột bên trái: ÉPISODE và số tập
            {\\[-40pt]
                \flaregothic\fontsize{20pt}{20pt}\selectfont \centering
                \color{eptype}HISTOIRE\color{black}\\
                \dawnland\fontsize{72pt}{72pt}\selectfont
                \color{epnum}3
                \\[18pt]
                \flaregothic\fontsize{14pt}{14pt}\selectfont
                \color{epattr}BIENVENUE, \uline{\color{black} Kaigou} !
                % \flaregothic\fontsize{18pt}{18pt}\selectfont
                % \color{epattr}ÉPISODE PERDU
                \color{black}
            }
             &
            % Dòng thứ nhất: tên tập tiếng Nhật
            {\fotskip\fontsize{18pt}{18pt}\selectfont \ruby{姉}{し}\ruby{妹}{まい}の\ruby{再}{さい}\ruby{会}{かい}！}  \\[6pt]
            % Dòng thứ hai: tên tập tiếng Anh
             & {\cafeteria\fontsize{18pt}{18pt}\selectfont Sister Reunion}                                \\[4pt]
            % Dòng thứ ba: tên tập tiếng Việt
             & {\vnmsans\fontsize{18pt}{18pt}\selectfont Đoàn tụ}                                         \\[8pt]
            % Dòng thứ tư: Nhân vật xuất hiện
             & {\caviar{\fontsize{14pt}{14pt}\selectfont Nhân vật: \emp{kuon1} \emp{kaigou1} \emp{suzuki1}}} \\[8pt]
        \end{tabular}
    \end{adjustwidth}
\end{table}
%===========================================NỘI DUNG==========================================
\color{black}

\pov{Hikawa Kuon}

Đã bao lâu rồi kể từ chúng tôi mắc kẹt ở nơi này nhỉ?

Mười phút? Ba mươi phút? Một tiếng? Tôi cũng chẳng biết nữa.

Tôi và cô bé tiếp tục\dots\ nói sao nhỉ, lê bước thì đúng hơn. Không còn là ``lang thang'' nữa, bởi cảm giác này giống như đang đi trong một giấc mơ không điểm kết. Mỗi bước chân đập xuống nền đất tối đen kia như thể không để lại dấu vết, không âm thanh, không trọng lượng. Thế giới này vẫn như cũ: kỳ quặc, tĩnh lặng và có phần siêu thực.

Chúng tôi lại bắt gặp một thứ kỳ lạ khác, một đường ống kim loại to bằng bắp chân, cắm chỏng chơ từ trên trời xuống như thể ai đó đâm nó từ tầng mây.

Một giọt nước nhỏ tí tách từ đầu ống rơi xuống, vỡ tan thành vệt sáng lấp lánh khi chạm đất. Tôi chưa kịp thắc mắc thì thấy Suzuki-chan khẽ giật bắn người, rồi siết chặt tay tôi một chút.

Tôi liếc mắt sang, thấy gương mặt cô bé hơi tái, bối rối quay đi. ``\vie{O-Ờm\dots\ Kuon-nii-san, c-chỉ là\dots}''

\dots Chắc tôi cũng đoán được lý do vì sao tại như thế. Một phần, bởi nơi này hoàn toàn vắng lặng, nên một tiếng động dù chỉ khe khẽ thôi cũng khiến cho cô bé giật mình.

Nhưng phần khác, tôi lại cảm thấy\dots\ cô bé có một điều gì đó rất khó nói, và tôi đoán được phần nào. Nhưng tôi không tiện nhắc, chỉ khẽ siết nhẹ lại như để nói: ``\vie{Ổn mà. Có anh ở đây rồi.}''

Một lúc sau, khi cả hai đã bình ổn lại nhịp thở và bước chân, tôi chợt nhớ ra điều gì đó: ``\vie{Này, Suzuki-chan?}''

``\vie{Dạ?}''

``\vie{Chị gái em trông như thế nào vậy? Anh đang nghĩ\dots\ ít nhất phải có một cái gì đó để anh có thể xác định được dễ dàng hơn,}'' tôi giải thích. Có lẽ vì tôi mải giúp em ấy dễ chịu hơn với tôi, mà bây giờ tôi mới nhớ tới điều này.

Suzuki-chan hơi bất ngờ trước câu hỏi đó, nhưng rồi, cô bé khoanh tay lại, gương mặt lộ rõ vẻ đang cố gắng hồi tưởng.

``\vie{Ưm\dots Onee-sama\dots\ mặc áo phông đen, trông giống kiểu của anh, nhưng form to hơn.}''

Tôi gật đầu. ``\vie{Gì nữa?}''

``\vie{Mặc\dots\ đầm đỏ dài tới gối. À, chị ấy buộc tóc đuôi ngựa, tóc hơi xoăn nhẹ\dots\ đi dép quai hậu giống như anh vậy đó.}''

Tôi lặng thinh vài giây rồi nhướng mày. ``\vie{Tức là\dots}''

Chưa kịp để tôi nói gì thêm, tôi để ý thấy cô bé cứ nhìn chằm chằm vào mặt tôi, soi xét từng chi tiết.

``\vie{Giờ em mới để ý, chị ấy nhìn hao hao giống anh đó, Kuon-nii-san ạ\dots}''

Tôi bật cười khẽ, một tiếng cười vừa mỉa mai vừa bất lực. ``\vie{Bây giờ em mới nhận ra à?}''

``\vie{Ơm\dots\ T-Tại lúc gặp anh, em lo quá nên đâu có để ý được gì đâu ạ\dots}'' Suzuki-chan tỏ vẻ bối rối, hai tay huơ lên huơ xuống như muốn bào chữa.

Câu nói ấy khiến tôi bật cười thật sự. Đúng là lúc ban đầu, cô bé như sắp khóc tới nơi. Vậy mà bây giờ đã có thể nói chuyện kiểu nhẹ nhàng thế này, nghĩ lại cũng đáng nể cho cả hai người chúng tôi thật.

Tôi trầm ngâm một chút, rồi quyết định nói thật. Tôi nghĩ nếu cứ giữ mãi như thế thì cũng không chắc tới khi nào tôi có thể nói ra.

``\vie{Suzuki-chan này\dots}''

``\vie{Dạ?}''

``\vie{Anh nghĩ\dots\ anh đã gặp chị em rồi.}''

Ngay lập tức, ánh mắt của Suzuki-chan mở to như hai cái đĩa. Cô ngập ngừng một chút rồi nắm lấy tay áo tôi, như thấy được một tia hy vọng nhỏ nhoi: ``\vie{T-Thật hả anh!? Chị ấy\dots\ Onee-sama đâu rồi!?}''

``\vie{Anh không biết,}'' tôi thở một hơi dài, rồi lắc đầu nhẹ, như không muốn dập tắt niềm hy vọng ấy. \vie{Chỉ nhớ lúc đó chị em- à không, chị gái đó tưởng anh là kẻ xấu, bảo là làm gì tệ với em. Cô ấy tấn công anh tới tấp. Sức mạnh thì như siêu nhân\dots}

Tôi vẫn còn nổi da gà vì cú đấm chết chóc lúc đó. Những cú đấm vừa nhanh vừa mạnh kia làm cho tôi gặp nhiều phen hú vía, có lúc làm tôi nghĩ rằng đã có thể bỏ mạng ở nơi này.

``\vie{Chị ấy đánh anh sao!?}'' cô bé hoảng loạn, nắm chặt tay áo của tôi hơn. ``\vie{Anh vẫn ổn chứ!?}''

``\vie{Cũng không hẳn là đánh, mà là suýt nữa nghiền anh ra bã.}''

``\vie{Ể!!??}'' Suzuki-chan xanh mặt, nhưng trước khi cô bé nói gì thêm, tôi vội chen lời: ``\vie{Nhưng mà anh ổn mà! Thật! Không có trọng thương gì đâu!}''

Phải để tôi cho ấy xem toàn thân của tôi vẫn còn toàn vẹn, cô bé mới có thể thở phào nhẹ nhõm.

``\vie{Cơ mà\dots}''

``\vie{Sao thế, Kuon-nii-san?}''

Tôi nhớ lại biểu cảm sau cùng của cô gái ấy: một khuôn mặt khắc khổ và bế tác đang gọi tên cô em gái của mình trong sự cùng cực, trái ngược hoàn toàn với sự dữ tợn và hùng hổ của một cô gái ghét đàn ông.

``\vie{Anh có nhớ rằng ngay khi anh đang ở thế lưỡng nan, cô ấy bỗng chốc dừng lại, rồi chợt bật khóc gọi tên em. Chắc là lúc ấy, chị em đang rất tuyệt vọng vì lạc mất em đấy.}''

Cô bé im lặng. Mắt cụp xuống. Gương mặt nhỏ nhắn dần co lại. Suzuki lí nhí, cảm thấy như có tội: ``\vie{\hl{E-Em xin lỗi\dots{} Em không nghĩ Onee-sama cố ý\dots}}''

Tôi vỗ vai cô bé, xoa nhẹ, trìu mến nhìn cô bé vừa quay trở lại tính cách rụt rè lúc đầu: ``\vie{Không cần xin lỗi đâu. Anh không trách chị em, cũng không trách em. Chỉ là, nếu đúng chị em là người đó, thì ít nhất anh có thể giúp hai người tìm lại nhau.}''

Cô bé lặng lẽ gật đầu, đôi vai run nhẹ, dường như thấu hiểu được lời tôi nói.

Tôi siết chặt bàn tay mình như thể đã tìm được con đường chân lý.

Chúng tôi không biết mình đang đi về đâu, nhưng chí ít có mục tiêu để đi. Và điều đó, trong thế giới vô định này, chính là thứ quan trọng nhất.

\pov{Hikawa Suzuki}

Tôi lặng thinh, mắt nhìn chằm chằm vào đôi giày vải đã sờn gót của mình. Trong tim tôi, những ký ức dồn về như một dòng suối lặng lẽ nhưng lạnh buốt.

Onee-sama lúc nào cũng vậy cả.

Từ khi tôi còn học tiểu học, mỗi lần tôi bị mấy đứa con trai trong lớp chọc ghẹo, đôi khi chỉ là trêu đùa nhẹ nhàng, đôi khi là ác ý, chị đều là người đầu tiên xuất hiện, kéo tôi ra sau lưng rồi nói bằng cái giọng lạnh tanh mà chỉ chị mới có: ``\vie{Xéo. Nếu không tao cho cắm bọn mày dính vào lưng luôn đấy.}''

Và lũ nhóc đó biến mất thật. Lúc nào cũng vậy.

Tôi biết Onee-sama chỉ muốn bảo vệ tôi. Tôi biết chị ghét đàn ông do đã từng bị tổn thương vì họ đã cười cợt, xúc phạm, và phản bội chị ấy.

Nhưng\dots\ nhiều lần tôi đã muốn nói rằng, không phải ai cũng có tâm địa xấu xa cả.

Có những người thực sự tốt, như thầy giáo của tôi hồi cấp hai, những người bạn học nam giúp tôi theo kịp bài học\dots

\dots hay bây giờ, như Kuon-nii-san, chẳng hạn.

Nhưng chị ấy có bao giờ nghe tôi nói. Chị không bao giờ cho tôi có cơ hội lên tiếng. Mỗi lần tôi định lên tiếng, chị sẽ siết tay tôi thật chặt, hoặc che tai tôi lại, như thể không muốn tôi bị ``làm hư'' bởi những lời đường mật của cánh đàn ông.

Tôi biết chị yêu tôi. Tôi biết.

Nhưng tình yêu ấy đôi khi khiến tôi nghẹt thở.

Tôi không trách chị, nhưng tôi sợ. Sợ một ngày nào đó, chị sẽ đẩy xa tất cả mọi người đến mức không còn ai muốn lại gần, kể cả tôi.

Tôi\dots\ tôi chỉ mong chị có thể bình tĩnh lại một chút thôi. Tôi chỉ muốn chị để tôi nói lên những lời thật lòng, dù chi là một chút.

Biết đâu, nếu Kuon-nii-san gặp lại chị ấy, anh có thể\dots\ giúp chị hiểu ra điều gì đó.

\dots

\dots

Tôi nuốt nước bọt khan. Trái tim vẫn còn đập nhanh, không hẳn vì sợ, mà vì lo. Lo cho Onee-sama, lo cho Kuon-nii-san, lo cho cả chính tôi.

Tôi thực sự không muốn anh ấy bị tổn thương chỉ vì tôi. Tôi muốn nói ra suy nghĩ của mình.

Và tôi nghĩ\dots\ tôi đã có đủ can đảm để làm điều này.

\il

Tôi khẽ xiết hai đùi lại.

Thành thật mà nói, cơ thể tôi đang gửi tín hiệu cầu cứu từ nãy giờ rồi. Cảm giác mỗi bước đi như đè thêm một lít nước nữa vào bàng quang đang căng cứng của tôi. Tôi cảm tưởng bây giờ nó đang to như cái trống rồi\dots

Tôi vẫn giữ vẻ bình tĩnh khi đi bên cạnh anh Kuon, vẫn mỉm cười, vẫn nói chuyện nhẹ nhàng\dots\ Nhưng chỉ khi anh ấy nhìn sang hướng khác, tôi mới dám lén giữ tay trước váy mình, ép chặt lại, mong sao giữ được lâu hơn một chút.

``\textit{Không sao cả. Chị từng dạy mình cách nhịn rồi. Mình chịu được mà\dots}'' tôi tự nhủ.

Nhưng\dots\ liệu mình còn chịu được bao lâu nữa đây\dots?

\il

\ct

\pov{Hikawa Kuon}

Một khoảng thời gian nữa lại trôi qua.

Tôi vẫn đang đi cùng Suzuki-chan, mắt dõi theo phía trước như mọi lần. The Void lúc này im ắng đến lạ thường, chỉ có tiếng bước chân nhè nhẹ của cả hai vang lên trên nền trống rỗng mơ hồ ấy.

Đầu tôi vẫn còn đang suy nghĩ tới người chị của cô bé.

Cô gái ấy, người đã tấn công tôi không chút do dự, lại có thể bật khóc gọi tên em gái mình chỉ trong khoảnh khắc, dù rằng đó có thể là một ảo giác. Sự đối lập ấy khiến tôi không khỏi băn khoăn: liệu đó là bản năng bảo vệ, hay là nỗi sợ mất đi người thân duy nhất?

Tôi còn đang định hỏi thêm vài điều về cô gái ấy thì bất chợt\dots

``\vie{\textbf{Suzu\dots\ Suzu\dots! Em ở đâu\dots?}}''

Một tiếng gọi từ xa vọng lại, méo mó và khản đặc. Nó không vang như tiếng hét, mà như một tiếng rên gằn trong cổ họng, vừa tuyệt vọng, vừa đầy lo lắng.

Nhưng chất giọng này không thể nhầm lẫn vào đâu được. Là \emph{cô ấy}. Tôi đứng khựng lại, nghiêng đầu.

``\vie{Anh có nghe thấy tiếng đó không\dots?}'' cô bé khẽ hỏi, hai mắt mở to, như vừa nhận ra một cái gì đó thân thương.

``\vie{\dots Có.}'' Tôi gật đầu, tim bắt đầu đập nhanh hơn.

Chính là cô gái mà tôi gặp lúc đầu. Là cô gái định hành tôi tới nhừ tử lúc trước. Chất giọng ấy, sự run rẩy trong từng chữ ấy, không thể nhầm vào đâu được.

Một bóng người phía trước xuất hiện đằng sau bóng tối. Một dáng hình quen thuộc, buộc tóc đuôi ngựa, vẫn mặc bộ đồ như lúc chạm mặt lần đầu, dù có phần lấm lem hơn, ánh mắt sục sạo giữa không gian vô định như tìm một điểm tựa.

Và rồi, ánh mắt xanh lam của cô ấy liếc sang tôi. Một thứ sát khí lại bừng bừng xuất hiện.

``\vie{\textbf{Lại là mày à!?}}'' cô rít lên như một con hổ bị thương, giọng chát đầy giận dữ.

Thấy rằng việc giải thích cho cô ấy là điều vô ích, tôi chỉ thở dài, nhắm mắt lại rồi lùi lại nửa bước, hít một hơi thật sâu.

``\eng{Ah sh*t, here we go again\dots}'' tôi lầm bầm. Vai nới lỏng, chân nhún nhẹ, chuẩn bị thế thủ để tránh đòn, vì tôi biết thể nào cũng sắp có một cú đấm từ địa ngục bay tới.

Nhưng ngay trước khi hai chúng tôi có thể làm được gì với nhau\dots

``\textbf{Onee-sama!!!}''

Tiếng hét của Suzuki-chan vang lên đầy xúc động, cắt ngang toàn bộ bầu không khí. Tôi quay đầu lại nhìn em ấy đang chạy tới, nước mắt lưng tròng, và rồi ôm lấy người chị thân thiết. Mọi thứ diễn ra quá nhanh, khiến không ai trong chúng tôi kịp phản ứng.

\dots Thôi thì, tôi quyết định nhường chỗ cảm xúc đoàn tụ nhỏ nhoi này cho hai chị em, còn tôi sẽ chỉ đứng nhìn từ xa.

\pov{Hikawa Suzuki}

Tôi không nghĩ ngợi gì nữa.

Ngay khi ánh mắt tôi bắt gặp chị ấy, dáng người cao ráo, tóc buộc đuôi ngựa, bước chân chậm chạp như đang vỡ ra dưới sức nặng của sự tuyệt vọng, tim tôi như ngừng đập một nhịp.

Mắt tôi mở to, cổ họng nghẹn ứ, hơi thở như tắc lại nơi lồng ngực. Tôi không tin nổi vào mắt mình. Là chị\dots\ chính là Onee-sama!

Tôi gào lên, không kịp suy nghĩ: ``\textbf{Onee-sama!!!}''

Cơ thể tôi tự động lao về phía trước, hai chân chạy như lướt trên không. Không gì có thể ngăn tôi lại lúc này, không phải nỗi sợ, không phải sự dè chừng, cũng chẳng phải cơn đau căng tức đang âm ỉ bên trong tôi từ nhiều giờ trước.

Onee-sama quay lại.

Onee-sama nhìn tôi.

Đôi mắt xanh lam ấy, ánh mắt luôn sắc lạnh, đề phòng với cả thế giới, giờ đây đang mở lớn đến sững sờ, như không tin vào mắt mình.

``Su\dots\ Suzu\dots?''

Chị lắp bắp như vừa nhìn thấy hồn ma, rồi run rẩy giơ tay ra, môi mấp máy: ``\vie{Em\dots\ em vẫn còn sống ư\dots!?}''

Tôi gật đầu lia lịa trong khi nước mắt bắt đầu trào ra. ``\vie{Vâng! Em đây! Là em mà, Suzuki đây!!}''

Chị khuỵu xuống ngay tại chỗ, đôi tay đưa ra, đón lấy tôi. ``\vie{Trời ơi\dots\ chị tưởng chị đã đánh mất em rồi chứ\dots}''

Tôi lao vào lòng Onee-sama, ôm chặt như thể sợ nếu buông ra thì chị sẽ tan biến mất. Cả người tôi rúc vào ngực chị, như ngày còn bé. Cả người chị run rẩy, từng hơi thở của chị nóng hổi táp vào mái đầu tôi.

``\vie{Em xin lỗi\dots\ em đã làm chị lo\dots}'' tôi vừa nức nở vừa nói, hai tay nắm chặt lấy tấm lưng quen thuộc đó.

``\vie{Không\dots\ là chị mới phải xin lỗi\dots\ chị đã không bảo vệ được em\dots}'' giọng chị khàn đặc, vừa khổ sở vừa nhẹ bẫng.

``\vie{Không đâu\dots\ chị đã làm mọi thứ rồi\dots\ em biết mà,}'' tôi khẽ đáp, lắc đầu trong vòng tay chị.

Chị siết tôi chặt hơn.

``\vie{Chị đã đi khắp nơi tìm em, gọi em mãi không thấy\dots\ Chị đã sợ\dots\ chị sợ em sẽ lại bị bọn đó làm hại\dots\ Chị thậm chí còn nghĩ\dots\ còn nghĩ rằng em đã\dots}'' Onee-sama nghẹn lại, câu nói không dứt, nhưng may mắn thay, nó đã không xảy ra.

Tôi ngẩng mặt lên, nhìn chị: ``\vie{Em không sao mà, chị thấy không? Em đứng trước mặt chị đây này\dots}''

``\vie{Chị biết chứ\dots\ nhưng\dots}'' Onee-sama đưa tay lên, khẽ vuốt lên má tôi, như muốn chắc chắn rằng tôi không phải là ảo ảnh, ``\vie{\dots chị vẫn không dám tin. Nhìn thấy em\dots\ mà chị cứ tưởng là mơ\dots}''

``\vie{Thế\dots\ chị muốn em véo chị một cái không?}'' tôi nửa đùa, nửa nghiêm túc, môi vẫn run vì xúc động. Cũng không quá ngạc nhiên với phản ứng của chị, đã quá lâu rồi chúng tôi mới gặp mặt nhau theo cách không ngờ tới, thậm chí tôi còn nghĩ đây vẫn chỉ là một giấc mơ\dots

Nhưng đây là một giấc mơ có thật. Onee-sama của tôi vẫn ở đó, tôi vẫn cảm nhận được hơi ấm từ người chị yêu quý của mình, và điều quan trọng nhất\dots\ chúng tôi vẫn còn sống.

Chị khựng lại một giây, rồi bật cười, một tiếng cười mỏng manh, run rẩy nhưng ngọt lịm như lần đầu tiên tôi nghe lại tiếng cười ấy sau bao năm.

``\vie{Suzu ngốc ạ\dots}''

Tôi cười theo, vùi mặt vào vai chị. Tất cả lo lắng, tủi thân, buồn tủi trong thời gian dài đằng đẵng ấy\dots\ cuối cùng cũng tan chảy trong vòng tay của người mà tôi luôn tin tưởng nhất.

``\vie{Chị hứa đi\dots\ từ giờ chị sẽ không đi đâu nữa mà không có em.}''

``\vie{Ừ, chị hứa. Llần sau, chị sẽ nắm tay em thật chặt\dots\ Chị sẽ không để em lạc mất nữa,}'' Onee-sama chị thì thầm như thề nguyền, tay siết tôi thêm lần nữa.

Tôi không nói gì thêm. Chỉ rúc vào lòng chị, cảm nhận được mùi quen thuộc từ áo chị, cảm nhận nhịp tim đang đập rộn ràng như tôi vậy.

Tôi biết, dù có chuyện gì xảy ra đi nữa, chúng tôi đã tìm lại được nhau ở một nơi xa lạ và đáng sợ này.

\pov{Hikawa Kuon}

``\vie{\textit{May quá rồi.}}''

Tôi đứng đó, giữ một khoảng cách nhất định. Hai chị em họ ôm chầm lấy nhau, hai hàng lệ cứ lã chã rơi, tay siết chặt lưng đối phương như thể không muốn rời lấy nhau một chút nào nữa. Tôi nghe loáng thoáng tiếng khóc, những lời xin lỗi lặp đi lặp lại, như muốn trút hết những khoảng thời gian không xác định lạc mất nhau vào khoảnh khắc ngắn ngủi này.

Tôi không chen vào. Đây vốn dĩ không phải là chỗ của tôi.

Tôi khẽ cúi đầu, liếc sang đồng hồ đeo tay vẫn đứng yên lúc 8:13, như thể thời gian không còn tồn tại ở đây. Tôi rút điện thoại ra kiểm tra, vẫn không có sóng, hiển nhiên. Mọi ứng dụng vẫn chạy.

Trong cuộc hội ngộ vỡ òa của hai người con gái ấy, tôi lặng lẽ quay gót, bước đi về phía đối diện, không để lại lời nào.

Mọi việc tôi cần làm cho họ, tôi đã làm rồi.

Tôi không nghĩ mình là người hùng. Chỉ là một cậu con trai bình thường, tưởng rằng trễ giờ bảo vệ đồ án mà phải bỏ dở công việc, để rồi lỡ bị cuốn vào cái xoáy lỗ đen quái quỷ, rồi tình cờ gặp một cô bé lạc lối, dẫn cô đến nơi cần đến. Tôi cũng chẳng cần một lời cảm ơn, càng không muốn ai phải mang ơn vì tôi cả.

Giờ thì xong xuôi cả rồi, họ đã đoàn tụ với nhau, tôi có thể tiếp tục tìm đường thoát ra khỏi cái chốn rỗng tuếch này. Một mình.

Mỗi bước chân như dẫm lên một thứ gì đó mềm nhũn vô hình, một cảm giác dẫu kỳ quặc lúc đầu, nhưng giờ tôi cũng quen rồi.

Cầm sẵn điện thoại trên tay, tôi mở ứng dụng phát nhạc và chọn album mà tôi định phát trước đó.

``\textit{Khúc Tiễn Biệt của Chuột Béo\dots\ Đây rồi.}''

Dẫu chỉ có âm nhạc để giúp tôi cầm cự ở chốn này và không còn ánh sáng chỉ đường, một mình tôi vẫn có thể tiếp tục cuộc hành trình này. Thật sự, chỉ muốn về nhà nhanh thôi.

Tôi dần dần tiến về màn đêm, bỏ lại bóng dáng hai người con gái đằng sau lưng.

``\textbf{Kuon-nii-san!}''

Tôi khựng lại. Tiếng của Suzuki-chan gọi tôi. Không một chút run rẩy. Không một chút dè dặt.

Tại sao chứ? Tôi đã hoàn thành nhiệm vụ rồi mà. Giờ tôi với cô bé có còn liên quan gì với nhau đâu. Càng không muốn dính líu gì tới cô gái hung hăng kia.

Tôi quay lại. Cô bé đã rời khỏi vòng tay của chị mình, chạy về phía tôi, tay nắm chặt gấu áo. Trên gương mặt nhỏ là sự bối rối pha lẫn quyết tâm.

``\vie{Kuon-onii-san\dots\ anh đừng đi. Anh\dots\ không phải người xấu. Em\dots\ em muốn Onee-sama tin anh.}''

Tôi nhìn cô bé, rồi nhìn chị gái của cô bé, người đang đứng đằng xa, đôi mắt vẫn nghi ngại nhưng không còn dữ tợn như lần đầu. Như thể cô ta đang chờ xem tôi sẽ làm gì.

Tôi thở dài, lẩm bẩm đầy mệt mỏi: ``\vie{Lại cái gì nữa đây\dots}''

``\vie{Em biết là anh không thích bị lôi vào rắc rối, nhưng\dots\ nếu anh đi, thì sẽ không ai có thể hiểu chị em cả.}''

Không để tôi nói thêm lới nào, Suzuki-chan đã kéo tôi lại gần đến cô gái kia mà tôi không hề hay biết.

\pov{Hikawa Suzuki}

Khi tôi quay đầu lại, tôi thấy Kuon-nii-san đang lặng lẽ rời đi.

Anh bước đi thật chậm, nhưng mỗi bước chân như dứt ra một phần gì đó trong tôi. Anh không quay lại. Không vẫy tay. Cũng chẳng để lại lời chào nào. Anh chỉ đơn giản là rời khỏi chúng tôi. Như thể anh nghĩ vai trò của mình đến đây là đủ.

Nhưng bản thân tôi lại không nghĩ vậy.

``Kuon-nii-san\dots!'' tôi khẽ gọi, rồi vội lao ra khỏi vòng tay chị mình.

``\vie{Em định làm gì vậy?}'' Onee-sama nắm lấy cổ tay tôi, giọng chị gắt lên. ``\vie{Em quan tâm tới hắn ta làm gì? Hắn ta đâu còn cần gì ở em nữa!}''

Tôi cắn môi. Onee-sama dường như vẫn không chịu hiểu, không chịu nghe những gì tôi nói. Tôi không thể để như thế này được thêm nữa.

``\vie{Không, không phải vậy đâu\dots}'' tôi nén lại cảm xúc, cố nói thật rõ. ``\vie{Anh ấy\dots\ chính anh ấy đã giúp em vượt qua khoảng thời gian lạc lõng ở nơi này. Chính anh ấy đã tìm thấy em, bảo vệ em, cho em đồ ăn, ở bên em suốt chặng đường đến tìm chị đó, Onee-sama à.}''

Chị nhìn tôi với đôi mắt nghi ngờ, rồi lầm bầm: ``\vie{Tên đó đã nhồi vào đầu em những thứ gì mà giờ em cứ khăng khăng bênh hắn như vậy\dots?}''

Tôi lắc đầu mạnh, giọng tuy nhỏ nhưng đầy dứt khoát: ``\vie{Không ai nhồi vào đầu em điều gì cả. Những gì em đang nói\dots\ là điều em thật sự nghĩ.}''

Tôi ngẩng lên nhìn thẳng vào đôi mắt chị, đôi mắt đã từng che chắn cho tôi khỏi biết bao nỗi sợ, nhưng giờ lại đang khép chặt vì sợ tổn thương.

``\vie{Em biết\dots\ em luôn là người để chị gồng gánh, bảo vệ. Nhưng em cũng muốn được bảo vệ chị, dù chỉ một lần thôi. Nếu chị không thể tin vào anh ấy\dots\ thì ít nhất, xin chị hãy tin vào em.}''

Tôi giương đôi mắt cún con lên mặt chị ấy, rồi khẽ nói bằng cả tấm lòng: ``\vie{Làm ơn đó, Onee-sama. Em cầu xin chị hãy cho em cơ hội giải thích.}''

Onee-sama nới lỏng tay ra. Gương mặt chị không nói gì, nhưng ánh mắt đã dịu lại.

Tôi không đợi thêm nữa.

Tôi quay người, chạy đi, mặc cho chị gọi với theo đằng sau. Tôi không chắc chị sẽ tha thứ cho anh ấy. Cũng không biết chị có hiểu hết lòng tôi hay không.

Nhưng tôi biết một điều: nếu không có Kuon-nii-san, tôi đã không thể nào tìm được chị ngay lúc này.

Tôi rướn bước, chạy về phía anh Kuon, tay giơ ra phía trước như muốn giữ lấy một điều gì đó đang sắp tuột khỏi tay. ``\vie{Kuon-nii-san, đừng đi mà!}''

Anh ấy khựng lại. Chậm rãi quay đầu, gương mặt hiện rõ vẻ bối rối xen lẫn mỏi mệt.

Tôi không nói gì, chỉ tiến đến gần và nắm lấy tay áo anh. ``\vie{Anh\dots\ đừng đi. Anh\dots\ không phải người xấu. Em\dots\ em muốn Onee-sama tin anh.}''

Anh nhìn tôi một lúc lâu, rồi khẽ thở dài, lẩm bẩm: ``\vie{Lại cái gì nữa đây\dots}''

Tôi cúi đầu, lí nhí: ``\vie{Em biết là anh không thích bị lôi vào rắc rối, nhưng\dots\ nếu anh đi, thì sẽ không ai có thể hiểu chị em cả.}''

Câu đó khiến anh ngẩn ra, còn tôi thì dắt tay anh quay ngược lại, không kịp để anh ấy nghĩ thêm điều gì. Đến khi quay về, Onee-sama vẫn đứng đó, hai tay khoanh trước ngực, ánh mắt trầm ngâm nhưng không còn lửa giận như trước.

``\vie{Đây là lần thứ hai chị chạm mặt hắn đấy, Suzu,}'' chị lên tiếng, giọng tuy lạnh lùng nhưng không còn sắc như dao.

``\vie{Em biết,}'' tôi đáp, bước lên một bước để đứng chắn giữa hai người họ. ``\vie{Nhưng em cũng biết là nếu không làm vậy, chị sẽ chẳng bao giờ chịu lắng nghe cả.}''

Chị im lặng, nhìn chằm chằm vào hai chúng tôi. Kuon-nii-san thì chỉ khẽ nhún vai, nhưng không bỏ đi nữa. Anh đứng đó, hai tay đút túi, mặt lộ vẻ nhẫn nại như thể chờ trận cuồng phong nổi lên bất cứ lúc nào.

Tôi hít một hơi thật sâu, rồi bắt đầu giãi bày những nỗi niềm từ tận đáy lòng mình.

``\vie{Onee-sama\dots\ Em biết chị ghét đàn ông. Em biết chị luôn nghĩ họ chỉ mang đến tổn thương. Quả thật, có những kẻ đã nhục mạ chúng ta, chê cười chúng ta là không có bố mẹ, hạ thấp chúng ta chỉ vì nghèo\dots}''

Cả hai anh chị vẫn im lặng. Tôi tiếp tục nói, giọng tôi không run, mà ngược lại, rõ ràng hơn bao giờ hết: ``\vie{Nhưng Onee-sama ạ, không phải ai cũng vậy. Kuon-nii-san là người đã giúp em, tôn trọng em, và không hề làm gì khiến em sợ cả.}''

Chị nheo mắt lại, hẳn nhiên càng tỏ vẻ nghi ngờ hơn: ``\vie{Chỉ vì hắn đối xử tốt với em một chút mà em đã nghĩ hắn là người tốt ư? Đầy người ở ngoài kia cũng giở giọng như thế, cuối cùng thì--}''

Kuon-nii-san vội cắt lời: ``\vie{Cứ để cho cô bé nói đi. Có gì thì cứ phản biện sau,}'' hai tay vẫn đút túi, tỏ vẻ bình tĩnh tới đáng sợ.

Onee-sama nghe vậy, nghiến răng tức lắm, nhưng cũng không làm gì được, đành im lặng tiếp tục nghe tôi nói.

``\vie{Em không tin anh ấy chỉ vì một hành động, mà vì\dots\ trong suốt thời gian ở cùng anh ấy, em chưa từng cảm thấy mình bị coi thường,}'' tôi lắc đầu, tiếp tục thuyết phục Onee-sama, với sự tự tin vốn có.

Tôi xoay người nhìn anh. ``\vie{Anh ấy không hề xem em là gánh nặng, không hề áp đặt điều gì lên em. Anh ấy chỉ đơn giản là luôn ở đó khi em cần. Anh ấy không đòi hỏi bất cứ điều gì từ em cả. Em biết chứ.}''

Tôi quay lại đối diện với chị.

``\vie{Em biết chị luôn muốn bảo vệ em. Nhưng\dots\ nếu chị bảo vệ em đến mức không cho em được tin ai khác, \emph{liệu em có còn được sống đúng nghĩa không?}}''

Chị trầm mặc, đôi môi khẽ mím lại. Tôi biết chị đang đấu tranh dữ dội trong lòng.

``\vie{\hl{Chị không cần phải tin anh ấy ngay lập tức,}}'' tôi nói nhỏ, thì thầm như sợ Kuon-nii-san nghe thấy, ``\vie{\hl{nhưng xin chị, ít nhất\dots{} hãy để anh ấy được bước đi bên em, cho đến khi em có thể về nhà cùng chị thật sự.}}''

Onee-sama nhìn tôi. Một cái nhìn kéo dài. Rồi chị chuyển ánh mắt sang anh ấy, người mới chỉ thở lấy một hơi dài.

Khoảnh khắc ấy, tôi chỉ dám nín thở mà chờ xem, liệu chị sẽ từ chối, hay cuối cùng, sẽ buông bớt tấm khiên trong lòng?

\pov{???}

Tôi nhìn đứa em gái bé nhỏ của mình, con bé vẫn đứng chắn giữa tôi và tên con trai kia như thể nó là người bảo vệ hắn vậy. Từng lời nó nói\dots\ từng âm tiết\dots\ đều thấm vào tôi một cách lặng lẽ mà không dễ chịu chút nào.

Tôi vẫn chưa mở miệng. Không phải vì tôi không có gì để nói, mà vì cổ họng tôi đang nghẹn lại.

Suzu bây giờ\dots\ không hề giống với con bé mà tôi biết. Từ lúc nào lại dám đứng lên nói như thế?

Từ bao giờ nó lại có thể nhìn thẳng vào mắt tôi mà phản đối lại lời của tôi vậy?

Tôi không biết, hay đúng hơn, tôi chưa từng nghĩ tới. Tôi chỉ biết là\dots\ mình thấy khó chịu. Không đơn thuần là vì nó cãi lời tôi nữa, mà đây là lần đầu tiên, tôi thấy mình không còn hiểu rõ em gái nữa.

Và rồi, hắn ta - cái tên tóc rối kia - mở miệng.

``\vie{Cô bé nói đúng đấy.}'' giọng hắn không to, không gắt, nhưng có một sự dứt khoát làm tôi giật mình.

Hắn đút tay túi quần, dáng điệu vẫn ung dung như cũ, nhưng ánh mắt thì không có vẻ xem thường tôi như những thằng đàn ông khác.

Ngược lại, chúng lại mang đến cho tôi một thứ gì khác\dots\ một sự cảm thông, nhưng cũng nghiêm khắc.

``\vie{Nếu cô cứ vơ đũa cả nắm, cho rằng tất cả đàn ông cô gặp đều là lũ rác rưởi\dots}'' hắn ngừng lại, mắt vẫn nhìn thẳng vào tôi, ``\vie{\dots thì cả đời cô sẽ chẳng thể tin ai được đâu.}''

Câu đó như một nhát chém xé toang lớp giáp mà tôi khoác lên trong lòng.

Sao hắn dám nói thẳng như thế? Mà không-- đúng hơn là\dots\ sao tôi không thấy ghét hắn khi hắn nói như thế?

Tôi cắn môi. Nhưng chưa kịp phản pháo gì, hắn đã nói tiếp: ``\vie{Tôi cũng không cướp em gái của cô đi để làm gì cả. Cũng chằng mong Suzuki-chan bênh vực tôi, cũng chẳng mong cô phải tin lời tôi nói.}''

``\vie{Và hơn hết,}'' hắn ta thở dài, tặc lưỡi, ``\vie{tôi cũng đếch muốn ở cái nơi chết dẫm này thêm một phút một giây nào nữa.}''

Lúc đó\dots\ có một điều gì đó vỡ ra trong lòng tôi.

Không phải vì lời lẽ mà hắn thốt ra. Chúng không tử tế. Chúng trần trụi, thô ráp, thiếu đi cái ngọt ngào mà tôi từng thấy ở những gã hay nịnh nọt.

Nhưng chính vì thế\dots\ tôi mới cảm thấy nó thật.

Tôi đã từng gặp biết bao thằng con trai, thằng đàn ông.

Đã từng nghe chúng xin lỗi, van vỉ, hứa hẹn đủ thứ mà chưa chắc chúng thực hiện được.

Từng thấy chúng dùng ánh mắt thương hại, bàn tay dối trá và giọng nói đường mật để thuyết phục hai chị em chúng tôi.

Nhưng chưa ai từng nói chuyện với tôi như cái cách mà hắn vừa nói.

Không rủ lòng thương. Không cố gắng làm tôi vừa lòng. Không cầu xin được tha thứ hay được tin tưởng.

Chỉ đơn giản là: ``Tôi như thế đấy. Tin thì tin, không tin thì thôi.''

Và lạ thật\dots\ Bỗng dưng, tôi cảm thấy bản thân mình muốn tin vào lời hắn nói.

Không hoàn toàn tin tưởng, tôi vẫn phải giữ cảnh giác. Nhưng ít nhất, đủ để tôi không gạt phăng hai đứa kia ra mà quay đi. Vì nếu tôi không tin tưởng hắn ta, chả phải là tôi đang phản bội Suzu sao?

Tôi nhìn cô em gái của tôi.  Nó đang nhìn tôi, ánh mắt van lơn, tay vẫn giữ lấy vạt áo tên đó.

Tôi buông một tiếng thở dài, tặc lưỡi một tiếng. ``\vie{Đi đâu thì đi cho nhanh đi,}'' tôi nói, cố giữ giọng thản nhiên. ``\vie{Cái nơi này tôi còn chán ngấy hơn cả các người.}''

\pov{Hikawa Kuon}

Tôi nghe thấy cô ấy nói: ``\vie{Đi đâu thì đi cho nhanh đi,}'' với cái giọng như thể đây là việc chẳng đáng để quan tâm. Nhưng sau tất cả những gì vừa xảy ra, tôi biết đó là lời nhượng bộ.

Tôi khẽ bật cười. Không phải là một tiếng cười vui vẻ gì cho cam. ``\vie{Cô nói là đi đâu thì đi, nhưng đến tôi còn chẳng biết đi đâu đây này\dots}'' Giọng tôi không giấu nổi sự chua chát.

Dẫu vậy, không hiểu sao\dots\ trong lòng tôi lại nhẹ đi một chút.

Không phải vì tôi thấy mình ``thắng''. Tôi chưa từng muốn phân thắng bại trong chuyện này.

Mà là vì\dots\ tôi thấy hy vọng.

Lần đầu tiên kể từ khi rơi vào nơi này, tôi có cảm giác như ba chúng tôi đã tạm thời cùng bước lên một điểm giao nhau.

Mọi thứ bắt đầu từ một cái lỗ xoáy, rồi một cái đệm, rồi một bữa ăn bị ăn vụng, rồi một cô bé đang rơi vào tuyệt vọng\dots

Và giờ, tôi đứng đây, đối diện với hai chị em vừa đoàn tụ, trong một cái cõi hư vô không bản đồ, không mục tiêu.

Tôi nhìn sang Suzuki-chan. Cô bé gần như nhảy bật vào lòng chị mình khi nghe thấy giọng nói có phần hậm hực kia của cô gái. Hai tay ôm chặt lấy người chị, nó dụi đầu vào vai chị, như thể sợ khoảnh khắc đó sẽ tan biến bất cứ lúc nào.

Tôi không thể không mỉm cười.

Ừm, đúng là cái ôm đó không dành cho tôi. Cái nhìn biết ơn kia cũng không hẳn hướng về phía tôi. Nhưng tôi biết, ít nhất một phần rất nhỏ trong chuyện này cái phần đã khiến cho cô ấy không quay lưng bỏ đi, đã khiến Suzuki-chan dám đứng ra bảo vệ tôi, đã khiến ba con người hoàn toàn xa lạ tạm thời ngừng đề phòng lẫn nhau - đều có một phần của tôi trong đó.

\dots Mà không, công lớn hơn chắc là cô bé rồi. Chính cô bé ấy là người đã thay đổi cục diện.

Cái giọng nhỏ nhẹ, đôi mắt ngân ngấn nước, và cả cái cách nó thốt lên: ``Nếu chị không tin anh ấy thì ít nhất\dots\ hãy tin vì em.'' Tôi nghĩ nếu là tôi, tôi cũng không nỡ quay lưng lại với lời nói đó.

Giờ cả hai người họ đang đứng trước mặt tôi. Một người ôm người kia. Còn tôi thì đứng hơi lệch sang bên, lưng vẫn cõng chiếc ba lô, trong túi thì đầy dây cáp và sách vở vô dụng, chẳng có cái gì thiết thực với tình huống hiện tại.

Và rồi, một khoảng im lặng lướt qua.

Tôi quyết định phá tan không khí bằng một câu (mà tôi cho là) vô hại: ``\vie{Ờm\dots\ nếu đã quyết định tạm đi chung rồi thì có lẽ giờ là lúc chúng ta tự giới thiệu một chút chứ nhỉ?}''

Tôi chưa biết tên của cô ấy là gì. Còn cô ấy thì chắc cũng chưa biết tên tôi.

Nhưng ít nhất, tôi cảm thấy cánh cửa đang dần mở ra, ngay giữa một nơi hoang vu hẻo lánh như thế này.

Tôi quay sang phía cô gái cao mấp mé bằng tôi.

Cô ấy, ờ, thôi thì giờ nhìn kỹ hơn rồi, tôi mới dám dùng từ xinh đẹp, tóc hơi xoăn, buộc đuôi ngựa lệch một bên, kiểu giống tôi nhưng nữ tính hơn. Gương mặt cân đối, mũi cao, mắt hai mí tự nhiên, được trang điểm nhẹ bằng phấn nền, eyeliner mảnh và son môi nhẹ màu cam đất. Tôi đoán là đồ trang điểm cũ được truyền lại từ mẹ của cô ấy.

Tôi thở nhẹ.

Và đúng lúc tôi đưa mắt nhìn xuống một chút, tôi chợt\dots\ nhận ra một điều nữa.

Là cái bộ phận đó. Cặp ngực. Dù không phải dạng quá khủng khiếp, vẫn khiến tôi phải giật mình. Phổng phao. Phải, to cỡ Suba-chan hoặc Mio-chan thật, và không phải theo kiểu độn, cũng không phải bơm silicon.

``\vie{Cậu nhìn vào cái gì của tôi đấy?}'' Giọng cô ấy sắc như dao cạo. Ánh mắt như mũi tên phóng tới. Tôi gần như vội ngẩng đầu nhìn thẳng lên trời.

``\vie{Không có gì hết,}'' tôi gãi đầu, cố vờ như đang gió thổi qua tai. ``\vie{Phân tích ngoại hình thôi.}''

``\vie{Phân tích cái con khỉ.}''

Suzuki-chan cười khúc khích, nhưng không dám phát ra tiếng.

Tôi ho khan một tiếng, vội chuyển sang chủ đề khác. ``\vie{À, tôi là Kuon, Hikawa Kuon. Đọc theo Hán Việt là Hỏa Xuyên Cửu Viễn. 25 tuổi. Sinh ngày 17 tháng Bảy.}''

``\vie{\dots Cái gì cơ?}'' cô ấy hơi rướn mày.

``\vie{Tên tôi, Hikawa Kuon,}'' tôi nhắc lại, rồi chợt nhíu mày khi cô ta vẫn giữ vẻ mặt khó hiểu. ``\vie{Có gì lạ sao?}''

Cô ấy hơi chậm rãi, rồi gật đầu, đôi mắt hơi trố ra như vừa phát hiện một cái gì đó mà cô chưa từng thấy. ``\vie{\dots Tôi cũng là Hikawa. Tên là Kaigou. \ruby{海}{かい}\ruby{強}{ごう}, \ruby{Cai}{Hải}-\ruby{gâu}{Cường}.}''

Tôi chớp mắt. ``Hikawa\dots\ Kaigou?''

Cô gật đầu, rồi đưa tay về phía chiếc túi nhỏ đang đeo bên hông, chắc định lấy điện thoại ra, nhưng lục tìm vài giây mà vẫn không thấy.

``\vie{Ủa\dots\ Điện thoại của mình đâu?}''

``\vie{Khoan đã\dots\ cái máy viền trầy trầy, đen đen đúng không?}'' tôi hỏi, lấy ra chiếc điện thoại đó.

``Ờm.''

Đúng cái máy đó, giờ đã đầy pin nhờ cục sạc dự phòng tôi mang theo. Phần góc hơi nứt kính một chút bởi đánh rơi.

``\vie{Của cô đây.}''

``\vie{Cái này\dots}'' Kaigou nhận lấy, ánh mắt hơi ngạc nhiên. ``\vie{Cậu lấy nó từ đâu?}''

``\vie{Từ sàn,}'' tôi nhún vai. ``\vie{Bị rơi ra lúc cô đấm tôi té chổng vó còn gì. Nứt một vết rồi đấy, cũng may là ở kính cường lực thôi.}''

Cô gái hơi cúi đầu. Không hẳn là có ý xin lỗi, nhưng ít nhất cũng không tỏ thái độ hằn học như trước.

Suzuki-chan lúc này đứng giữa hai chúng tôi, mắt nhìn bên này, rồi lại bên kia, như đang soi gương. Mà không, ngẫm nghĩ lại, tôi có cảm tưởng rằng người đang đối diện mình là một bản sao của mình vậy.

``\vie{Ủa? Cả hai người đều 25 tuổi, sinh cùng ngày\dots\ lại còn họ Hikawa nữa ạ?}''

Tôi và Kaigou đồng loạt nhìn nhau.

``\vie{Đừng nói là\dots}'' tôi lên tiếng trước, cô ấy phản xạ ngay sau đó: ``\vie{Không thể nào\dots}''

Suzuki chắp tay sau lưng, nở nụ cười tinh nghịch: ``\vie{Anh chị đúng là\dots\ sinh đôi á?}''

``\vie{\textbf{Không!}}'' cả hai chúng tôi cùng hét lên, dù rằng có hơi hoảng.

Dù vậy, trong bụng tôi vẫn không tránh khỏi cảm giác kỳ quặc. Hai người xa lạ, từ hai thế giới khác nhau, nhưng lại cùng tuổi, cùng sinh nhật, cùng họ, theo cả nghĩa cách đọc và cách viết\dots

Càng nghĩ, tôi càng thấy như có một thế lực vô hình nào đó cố tình sắp xếp ra cái tình huống tréo ngoe này vậy.

Trong lúc chúng tôi vẫn chưa khỏi bàng hoàng, cô bé khẽ giơ tay lên, phá tan bầu không khí khó xử giữa hai chúng tôi: ``\vie{Còn em là Hikawa Suzuki (\ruby{火}{ひ}\ruby{川}{かわ}\ruby{鈴}{すず}\ruby{木}{き}, \ruby{Hi}{Hỏa}-\ruby{ca-oa}{Xuyên}\ \ruby{Xư-dư}{Linh}-\ruby{ki}{Mộc})! 15 tuổi, học lớp 8, vì\dots\ em học trễ một chút.}''

Tôi nhìn con bé. Nó cười bẽn lẽn, tỏ vẻ hơi ngượng ngùng khi tự khai ra như thế, tay đung đưa nhẹ trước bụng, đôi mắt ngước nhìn tôi, rồi nhìn sang chị nó.

Ba người - ba họ Hikawa - từ hai vũ trụ khác nhau. Không biết đây là định mệnh, hay chỉ là trò đùa của The Void mang tới.

Nhưng tôi chắc chắn một điều: chuyến hành trình này mới chỉ là khởi đầu.

\il

Tôi vẫn chưa thật sự hết căng thẳng sau màn đối đáp với cô chị khó ưa kia, nhưng ít nhất thì chuyện cũng tạm yên. Suzuki-chan trông như trút được gánh nặng. Cô bé thở phào nhẹ nhõm, hai vai xụ xuống.

Tiếc rằng, sự nhẹ nhõm đó chỉ kéo dài được đúng năm giây.

``A\dots\ Aaah!!''

Tôi nghe tiếng hét nhòe trong cuống họng, rồi thấy cô bé cuống cuồng, hai tay chụm lại thật chặt phần bụng dưới, mắt nhìn quanh hoảng loạn như thể vừa nhận ra một điều cực kỳ khẩn cấp.

\dots Tôi còn nhìn thấy vài giọt nước vàng đang long tong từ dưới váy, nhưng tôi không lên tiếng. Khá thô thiển nếu tôi nói vậy.

Kaigou quay ngoắt sang, nét mặt hiện rõ vẻ lo lắng: ``\vie{Suzu!? Chuyện gì vậy!?}''

``\vie{Em\dots\ em\dots\ em phải đi tiểu\dots}'' Suzuki-chan líu cả lưỡi, gần như suýt khóc.

``\vie{\textbf{Cái gì cơ!? Bây giờ á!?}}'' cô chị gái càng hoảng hốt hơn. ``\vie{Không, cố nhịn thêm chút đi, để chị tìm nhà vệ sinh!}''

Tôi nghĩ ở một nơi mà phương hướng không rõ ràng này, tìm được một chỗ để ``giải quyết nỗi buồn'' một cách tử tế đúng nghĩa đã là điều bất khả thi rồi.

Tôi thở dài, nụ cười méo xệch nở ra trên mặt. Tôi đã nghi ngờ từ sớm rồi. Khi cô bé đi chậm, khi từ chối tôi cõng, và khi ánh mắt len lén quay đi mỗi lần tôi hỏi gì đó riêng tư. Mãi cho tới khi cái cách mà Suzuki-chan giật mình nhăn mặt khi nghe thấy tiếng nước nhỏ, mọi thứ mới được làm sáng tỏ.

``\vie{Cô không nhận ra à?}'' tôi quay sang Kaigou. ``\vie{Tôi đoán ra từ lúc trước rồi. Con bé chịu đựng cả quãng đường, và có khi là từ trước khi gặp tôi đấy.}''

Kaigou nghiến răng, vẻ bất lực hiện rõ. ``\vie{Nhưng\dots\ nhưng ở đây có nhà vệ sinh đâu!}''

``\vie{Ồ, không phải không có.}''

Tôi chỉ tay về một hướng hơi khuất, nơi có mấy phiến đá nhô lên, đủ để che tầm mắt từ chỗ chúng tôi. ``\vie{\emph{Chỗ đó.} Xa vừa đủ để ta có thể canh chừng.}''

Cô chị trố mắt lên, rõ ràng phản đối kịch liệt. ``\vie{Cái gì!? Cậu định để em gái tôi--}''

``\eng{Look,}'' tôi giơ bàn tay lên cao, vội ngắt lời cô. ``\eng{When you gotta go, you gotta go, wherever the place is.}\ (Mắc quá thì cứ đi thôi, bất kể chỗ đó là gì.) \vie{Cô thực sự muốn nhìn em gái cô tè dầm tại chỗ à?}''

Cô chị còn định cãi nữa, nhưng tiếng cô em gái rên khe khẽ phía sau khiến mọi lý lẽ dừng lại. Không chờ thêm giây phút nào, Suzuki-chan gật đầu lia lịa, rồi cắm đầu chạy về phía tôi vừa chỉ, bước chân khập khiễng như đang cố giữ cả một biển cả bên trong bụng. Nhìn mà tôi chỉ muốn rùng mình thay.

``\textbf{\vie{Khoan đã, Suzuki-chan!}}''

Cô bé khựng lại giữa chừng, quay đầu lại nhìn tôi đầy lo lắng, khuôn mặt đau khổ vì ``cánh cửa xả lũ'' đã gần ngay trước mắt, chỉ để bị tôi cản lại.

``\vie{C-Cái gì nữa\dots\ Kuon-nii-san?}''

``Đây.'' Tôi đưa một gói khăn giấy nhỏ, ném nhẹ về phía cô bé.

``Áaaa\dots!'' Suzuki-chan rít lên một tiếng đầy bối rối, mặt đỏ như cà chua, rồi vồ lấy gói giấy và phóng mất dạng sau mấy phiến đá.

Tôi bật cười nhẹ, rồi quay lưng lại, đeo tai nghe lên, chỉnh âm lượng vừa đủ. Vừa đủ để át đi những gì không nên nghe. Cũng tiện thể\dots\ tôi nhắm mắt lại luôn, gọi là ``đóng giả mù giả điếc trọn gói'', hoặc ít nhất là không bị cô chị coi tôi là kẻ biến thái.

``\vie{Cậu đang làm gì thế?}'' giọng Kaigou vẫn đầy hoài nghi. ``\vie{Sao lại đưa giấy lau cho Suzu?}''

``\vie{Đừng nói với tôi là cô thật sự muốn Suzuki-chan đi lại với phần dưới ướt sũng đấy nhá?}'' tôi đáp không quay lại, một bên tai nghe vẫn lủng lẳng, trước khi đeo hoàn toàn.

Kaigou không nói thêm gì nữa. Mà tôi nghĩ cũng chẳng cần cô ấy phải nói. Đó thực sự là một câu hỏi ngu.

Tôi thả mình vào tiếng nhạc, đếm từng giây trôi qua, chờ cho gió mang đi hết sự căng thẳng cuối cùng.

\pov{Hikawa Kaigou}

Tôi cứ đứng đó, tay khoanh lại, mắt liếc sang cô em gái đang giải quyết chuyện cá nhân, lúc lại liếc nhìn cậu ta - Hikawa Kuon.

Lẽ ra, trong tình huống như thế này, bất kỳ thằng đàn ông nào cũng sẽ giở trò. Lợi dụng lúc con bé yếu thế, mất cảnh giác mà giở trò đồi bại, hoặc chí ít cũng liếc trộm. Tôi đã thấy nhiều rồi. Quá nhiều là đằng khác.

Nhưng Kuon\dots\ cậu ta không làm gì cả.

Chỉ đơn giản là đưa cho Suzu gói giấy, một hành động đầy tế nhị mà cậu ta không bắt buộc phải làm, rồi quay lưng lại, đeo tai nghe vào, nhắm mắt như thể cả thế giới chẳng còn liên quan gì tới mình. Không nhìn, không nghe, không can dự vào.

Tôi nhìn bóng lưng cậu ta một lúc, lòng có chút chao đảo.

\dots Cậu ta không giống lũ còn lại.

Không lẽ\dots\ tôi đã sai?

Tôi khẽ rùng mình. Một phần bởi sự khác biệt rõ ràng giữa Kuon và đám đàn ông hạ đẳng kia là một trời một vực. Tôi chưa bao giờ gặp một người con trai nào lại tốt tính như thế.

Một phần khác\dots

``\textit{\vie{Mình cũng muốn đi tè nữa\dots\ Ư\dots\ mót quá\dots}}'' tôi khẽ lầm bầm, hai chân đan lại với nhau.

Từ lúc tôi và Suzu bị mắc kẹt ở nơi này vào một khoảng thời gian quá dài (mà tôi không thèm để tâm), vì quá mải đi tìm em mà đã không nghĩ tới chuyện cá nhân đó.

Cảm giác của con bé khi được ``giải phóng'' lúc này hẳn phải thấy sung sướng lắm.

Nhưng mà không được!

Tôi cắn răng, siết chặt hai tay ở bụng dưới lại.

Tôi là chị lớn. Tôi không thể làm điều đó ở đây. Nhất là trước mặt Suzu, còn chưa kể đến trước mặt cậu ta nữa chứ! Biết đâu cậu ta sẽ lại giở bài gì hại tôi thì sao?

Tôi vẫn còn đủ sức để nhịn được.

\pov{Hikawa Suzuki}

Sau một hồi giải quyết nhu cầu cá nhân dài tới mức tôi còn không biết được, tôi bẽn lẽn trở lại, từng bước nhỏ khẽ chạm đất, đôi gò má nóng rực như muốn bốc khói.

``\vie{C-Cảm ơn anh, Kuon-nii-san\dots}'' tôi lí nhí nói, tay siết vào vạt áo, không dám nhìn thẳng vào mắt anh. Sau cú lỡ miệng thốt lên đó, tôi chẳng biết liệu phải tìm chỗ nào để trốn đây\dots

Anh ấy quay sang, nhìn vũng nước mà tôi vừa thải ra. Một nụ cười tinh nghịch nhếch lên. ``\vie{Nếu Olympic có tổ chức bộ môn nhịn tiểu thì chắc em ẵm luôn huy chương vàng rồi đấy. Cái vũng như ổ voi kia chắc phải tầm ba lít rưỡi, bốn lít là ít!}''

Tôi suýt hét lên vì xấu hổ. ``K-Kuon-nii\dots!!'' tôi lấy tay che mặt, còn mặt thì nóng hơn cả nồi lẩu. ``\vie{Đừng có nhắc đến cái đó chứ\dots}''

*BỘP!* ``Owie.''

Tiếng va nhẹ vang lên. Tôi hé mắt ra thì thấy Onee-sama đang cụng nhẹ đầu anh ấy bằng trán. Rõ ràng chị ấy không thể dùng tay chân được, nên đó là cách nhẹ nhất để ``trừng phạt'' anh ấy rồi.

``\vie{Đừng có nói như thế!}'' chị lườm Kuon-nii-san, nhưng chẳng hiểu sao trong ánh mắt chị không có sự tức giận. Chỉ là\dots\ một chút chán nản pha lẫn ngạc nhiên.

``\vie{Đùa tí thôi mà, sao căng thế,}'' anh ấy xoa đầu, bật cười nhẹ.

Tôi khẽ cười. Họ đúng là lạ thật.

Sau đó, tôi quay sang anh ấy, rụt rè nhờ một việc mà bây giờ tôi mới có can đảm để nói: ``A-A-nô\dots\ \vie{Kuon-nii-san, nếu được\dots\ em có thể\dots\ được cõng không ạ?}''

Anh nhìn tôi, thoáng một vẻ bất ngờ. Chắc vì lúc trước tôi đã từ chối không chỉ một, mà vài ba lần rồi.

``Em nói cái gì thế Suzu!?'' Onee-sama cũng ngỡ ngàng. Cũng phải thôi, đây là lần đầu tiên tôi nhờ một người nào đó ngoài chị ấy cõng tôi lên mà.

Tôi siết vạt áo, rồi ngẩng đầu lên, đôi mắt ánh lên sự tự tin, nhìn thẳng vào hai người. ``\vie{Vì em\dots\ tin anh ấy ạ.}''

Cả hai người kia đều sững lại. Kuon thì nhìn tôi, đôi mắt ánh lên sự dịu dàng không gượng gạo. ``\vie{Được chứ.}''

Còn chị tôi thì không nói gì, cũng không hề lên tiếng phản đối.

Tôi bước lại gần, nhẹ nhàng trèo lên lưng anh ấy.

Lần đầu tiên kể từ khi rời khỏi thế giới của mình, tôi cảm thấy an tâm đến thế.

Có khi nào\dots\ tôi có thể coi anh ấy là anh trai của tôi chăng?

\end{document}